% Generated draft: requires mathematical and source review.
\documentclass[11pt]{article}
\usepackage[margin=25mm]{geometry}
\usepackage{fontspec}
\setmainfont{Latin Modern Roman}
\usepackage{amsmath,amssymb,mathtools,unicode-math}
\setmathfont{Latin Modern Math}
\usepackage{enumitem,needspace,xurl,hyperref}
\setlength{\parindent}{0pt}
\setlength{\parskip}{5pt}
\setlength{\emergencystretch}{3em}
\newcommand{\N}{\mathbb{N}}
\newcommand{\Z}{\mathbb{Z}}
\newcommand{\Q}{\mathbb{Q}}
\newcommand{\R}{\mathbb{R}}
\newcommand{\C}{\mathbb{C}}
\newcommand{\F}{\mathbb{F}}
\newcommand{\A}{\mathbb{A}}
\newcommand{\PP}{\mathbb{P}}
\newcommand{\E}{\mathbb{E}}
\newcommand{\eps}{\varepsilon}
\newcommand{\dd}{\,\mathrm{d}}
\begin{document}
\section*{UnsolvedMath ID: 30002312}
\textbf{Source problem code:} catalogue entry
\subsection{Definitions and mathematical statement}
\textbf{Recorded question.} \emph{Why is homogenization of the nonlinear plasticity problem difficult, particularly for its differential inclusion and the limited regularity of two-scale solutions?}

\textbf{Formulation limitation.} This is an explanatory question, rather than a specified mathematical conjecture. The supplied record gives no governing equations, differential inclusion, constitutive assumptions, unknowns, spatial or temporal domain, boundary or initial conditions, microscopic scaling, periodic or stochastic framework, solution spaces, convergence topology, or proposed homogenized problem. In particular, neither the relevant two-scale solutions nor their asserted regularity limitations are defined. Therefore no precise quantified homogenization statement can be recovered without additional source information.

The supplied statement-verification note labels the statement \texttt{unrecoverable}. It reports that the official workshop report immediately answers the explanatory question by identifying nonlinear equations, a differential inclusion, and insufficient regularity, and then describes a finite-element workaround. This is supplied source-verification information, not an independently checked account. No separate mathematical conjecture or claim concerning that workaround is specified here.
\subsection{Short English statement}
The record asks why nonlinear plasticity is difficult to homogenize, emphasizing a differential inclusion and limited two-scale regularity. It does not specify a mathematical model or a theorem to establish. A precise problem statement requires further source information; the supplied partially solved label does not resolve this formulation gap.
\subsection{Sources}
\begin{itemize}
\item Local UnsolvedMath registry, permanent ID \texttt{30002312}, title \emph{Homogenization of Nonlinear Plasticity}, source problem code \texttt{catalogue entry}, in the Oberwolfach Reports Open Problems collection. The recorded year is 2013; no individual proposer is supplied. The registry retains the original explanatory question for provenance and labels its statement status \texttt{unrecoverable}.
\item Supplied report citation: \emph{Interplay of Theory and Numerics for Deterministic and Stochastic Homogenization} (2013), Oberwolfach Reports, DOI 10.4171/owr/2013/14, pp.\ 801--865. Supplied source URL and literature-source link, labeled Oberwolfach Report 14/2013: \url{https://doi.org/10.4171/owr/2013/14}. The collection describes these reports as documenting mathematical workshops and their open-problem discussions.
\item Additional supplied literature source: arXiv:1604.02291, \url{https://arxiv.org/abs/1604.02291}. The record supplies no title, authors, or pinpoint result for this reference.
\item Literature assessment dated exactly \texttt{2026-08-21}: the supplied curation labels the topic partially solved and states that rigorous homogenization has been proved for important plasticity models, including stochastic hardening systems, while a nonlinear differential inclusion and insufficient two-scale regularity continue to obstruct a general theory. This is supplied dated curation derived from the OWR material, not independently verified literature assessment. Its time-sensitive status and cited claims require independent verification.
\end{itemize}
\end{document}
